\section{Conclusion}
\label{sec:conclusion}

This paper has studied authority allocation in a firm that relies on an external agentic-AI platform for task execution. The central argument is that the organizational consequences of generative AI cannot be understood by looking only at current task-level efficiency. Because delegation to external AI affects the firm's retained human recovery capability, the relevant organizational trade-off is intertemporal: greater current execution efficiency may come at the cost of weaker future resilience.

The analysis delivers three main conclusions. First, the optimal organization generally takes a hybrid form in which routine tasks are delegated to AI while sufficiently exceptional tasks remain under human authority. Second, full automation may be suboptimal even when AI dominates human handling task by task in normal operations, because authority today helps preserve internal recovery capability for future contingencies. Third, platform pricing distorts the firm's internal authority structure relative to a technological benchmark, affecting not only whether AI is adopted but also how authority is allocated conditional on adoption, with stronger distortions in longer-horizon execution environments.

More broadly, the paper suggests that the economics of generative AI is not only about the substitution of machines for labor. It is also about how dependence on externally supplied AI reshapes internal governance, organizational learning, and the boundaries of human discretion. In that sense, the relevant question is not simply how much AI a firm should use, but how a firm should organize itself when critical execution capabilities are supplied from outside the organization.
